\documentclass{article}


\usepackage[margin=1in]{geometry}
\usepackage{amsmath,amssymb}
\usepackage{natbib}

\providecommand{\keywords}[1]{\vspace{0.5em}\noindent\textbf{Keywords:} #1}
\providecommand{\And}{\and}

\usepackage[utf8]{inputenc} % allow utf-8 input
\usepackage[T1]{fontenc}    % use 8-bit T1 fonts
\usepackage{hyperref}       % hyperlinks
\usepackage{url}            % simple URL typesetting
\usepackage{booktabs}       % professional-quality tables
\usepackage{amsfonts}       % blackboard math symbols
\usepackage{nicefrac}       % compact symbols for 1/2, etc.
\usepackage{microtype}      % microtypography
\usepackage{lipsum}
\usepackage{fancyhdr}       % header
\usepackage{graphicx}       % graphics
\graphicspath{{media/}}     % organize your images and other figures under media/ folder

%Header
\pagestyle{fancy}
\thispagestyle{empty}
\rhead{ \textit{ }} 

% Update your Headers here
\fancyhead[LO]{MARL for Adversarial Spatial Decision-Making}

%% Title
\title{Multi-Agent Reinforcement Learning for Adversarial Spatial Decision-Making with Temporal Constraints}

\author{
  \begin{minipage}{0.45\textwidth}
    \centering
    Chaitanya Chaudhary \\
    Dept. of Computer Science \& Engineering \\
    IIT Gandhinagar \\
    \texttt{chaitanya.chaudhary@iitgn.ac.in}
  \end{minipage}%
  \begin{minipage}{0.45\textwidth}
    \centering
    Dhruv Meshram \\
    Dept. of Computer Science \& Engineering \\
    IIT Gandhinagar \\
    \texttt{dhruv.meshram@iitgn.ac.in}
  \end{minipage}
  \\[3ex]
  \begin{minipage}{0.45\textwidth}
    \centering
    Abhishek Tagalpallewar \\
    Dept. of Mechanical Engineering \\
    IIT Gandhinagar \\
    \texttt{abhishek.tagalpallewar@iitgn.ac.in}
  \end{minipage}%
  \begin{minipage}{0.45\textwidth}
    \centering
    Vedant Bondekar \\
    Dept. of Computer Science \& Engineering \\
    IIT Gandhinagar \\
    \texttt{vedant.bondekar@iitgn.ac.in}
  \end{minipage}
}


\begin{document}
\maketitle


\begin{abstract}
Multi-agent reinforcement learning (MARL) studies sequential decision-making problems in which multiple learning agents interact within a shared environment, making strategic adaptation central to the learning process. In competitive settings, the behavior of each agent depends not only on the environment but also on the evolving policy of an opposing agent, creating a natural connection between reinforcement learning and game-theoretic decision-making. In this work, we investigate these dynamics using \textbf{Modified Wallz}, a deterministic, sequential, two-player zero-sum environment in which agents simultaneously pursue opposing objectives while modifying the spatial constraints of the board. The environment extends the original wall-based mechanics with reusable walls, a one-time wall-bypassing hop, and temporary cell blocking, introducing resource constraints, temporal dependencies, and dynamically changing action availability. These mechanisms provide a controlled setting for studying how strategic resources and evolving constraints affect multi-agent policy learning. We formulate the environment in terms of agent observations, actions, rewards, and state transitions, and propose to evaluate MARL agents against heuristic and search-based baselines. For smaller configurations, exhaustive game-tree analysis and minimax-based reasoning will be investigated as references for learned behavior and computational feasibility. By varying resource and temporal parameters, we aim to characterize their effects on strategic adaptation, learning performance, and computational complexity, thereby providing insight into multi-agent decision-making under dynamic spatial and resource constraints.
\end{abstract}


% keywords can be removed
\keywords{First keyword \and Second keyword \and More}


\section{Introduction}

Multi-Agent Reinforcement Learning (MARL) considers sequential decision-making problems in which multiple adaptive agents interact within a shared environment. Such settings arise whenever the outcome of an agent's actions depends on the behavior of other decision-makers, making learning fundamentally different from conventional single-agent reinforcement learning. Markov games provide a natural mathematical framework for these settings by allowing multiple agents with interacting or competing objectives to influence a common environment \cite{littman1994markov}. MARL has consequently developed strong connections with game theory, optimization, and reinforcement learning, with applications spanning cooperative, competitive, and mixed-motive environments \cite{busoniu2008comprehensive,zhang2020multiagent}.

A central difficulty in MARL is that agents learn simultaneously. From the perspective of an individual learner, the effective environment can therefore become non-stationary as the policies of other agents change during training \cite{gronauer2022multiagent,zhang2020multiagent}. In addition, reasoning over multiple agents introduces joint action spaces whose size grows with the number of decision-makers, increasing the computational burden of policy search. These issues are particularly relevant in competitive settings, where an agent must learn not only how to maximize its own objective but also how its decisions affect, and are affected by, an adaptive opponent. Consequently, strategic adaptation rather than isolated action optimization becomes an important component of the learning problem.

We investigate these issues using \textbf{Modified Wallz}, a deterministic, sequential, two-player, zero-sum board environment. Each player attempts to reach the opponent's side while moving a pawn and placing walls that restrict movement subject to path-preservation constraints. The modified environment introduces reusable walls, a one-time opportunity to hop over an opponent's wall under defined legal conditions, and temporary cell blocking for a fixed number of future moves. These mechanisms introduce resource constraints and temporal dependencies into a fully observable competitive environment, allowing strategic interaction to be studied while keeping the underlying state and transition structure explicit.

The resulting environment provides a controlled setting for investigating how adaptive opponents, reusable strategic resources, and dynamically changing spatial constraints influence policy learning. We formulate the environment as a multi-agent sequential decision problem and aim to evaluate MARL agents against heuristic and search-based baselines. For smaller configurations, exhaustive game-tree analysis and minimax reasoning will also be considered to provide references for learned policies and to examine the tractability of exact computation. By varying the resource and temporal parameters, we seek to characterize their effects on strategic behavior, learning performance, and computational complexity. This motivates the formal definitions of the environment, agents, actions, rewards, and state transitions developed in the following sections.


\section{Headings: first level}
\label{sec:headings}

\lipsum[4] See Section \ref{sec:headings}.

\subsection{Headings: second level}
\lipsum[5]
\begin{equation}
\xi _{ij}(t)=P(x_{t}=i,x_{t+1}=j|y,v,w;\theta)= {\frac {\alpha _{i}(t)a^{w_t}_{ij}\beta _{j}(t+1)b^{v_{t+1}}_{j}(y_{t+1})}{\sum _{i=1}^{N} \sum _{j=1}^{N} \alpha _{i}(t)a^{w_t}_{ij}\beta _{j}(t+1)b^{v_{t+1}}_{j}(y_{t+1})}}
\end{equation}

\subsubsection{Headings: third level}
\lipsum[6]

\paragraph{Paragraph}
\lipsum[7]

\section{Examples of citations, figures, tables, references}
\label{sec:others}
\lipsum[8] \cite{kour2014real,kour2014fast} and see \cite{hadash2018estimate}.

The documentation for \verb+natbib+ may be found at
\begin{center}
  \url{http://mirrors.ctan.org/macros/latex/contrib/natbib/natnotes.pdf}
\end{center}
Of note is the command \verb+\citet+, which produces citations
appropriate for use in inline text.  For example,
\begin{verbatim}
   \citet{hasselmo} investigated\dots
\end{verbatim}
produces
\begin{quote}
  Hasselmo, et al.\ (1995) investigated\dots
\end{quote}

\begin{center}
  \url{https://www.ctan.org/pkg/booktabs}
\end{center}


\subsection{Figures}
\lipsum[10] 
See Figure \ref{fig:fig1}. Here is how you add footnotes. \footnote{Sample of the first footnote.}
\lipsum[11] 

\begin{figure}
  \centering
  \fbox{\rule[-.5cm]{4cm}{4cm} \rule[-.5cm]{4cm}{0cm}}
  \caption{Sample figure caption.}
  \label{fig:fig1}
\end{figure}

\subsection{Tables}
\lipsum[12]
See awesome Table~\ref{tab:table}.

\begin{table}
 \caption{Sample table title}
  \centering
  \begin{tabular}{lll}
    \toprule
    \multicolumn{2}{c}{Part}                   \\
    \cmidrule(r){1-2}
    Name     & Description     & Size ($\mu$m) \\
    \midrule
    Dendrite & Input terminal  & $\sim$100     \\
    Axon     & Output terminal & $\sim$10      \\
    Soma     & Cell body       & up to $10^6$  \\
    \bottomrule
  \end{tabular}
  \label{tab:table}
\end{table}

\subsection{Lists}
\begin{itemize}
\item Lorem ipsum dolor sit amet
\item consectetur adipiscing elit. 
\item Aliquam dignissim blandit est, in dictum tortor gravida eget. In ac rutrum magna.
\end{itemize}


\section{Conclusion}
Your conclusion here

\section*{Acknowledgments}
This was was supported in part by......

%Bibliography
\bibliographystyle{unsrt}  
\bibliography{references}  


\end{document}
